%%% Класс документа
\documentclass[a4paper,14pt]{article}

%%% Работа с русским языком
\usepackage{cmap}					% поиск в PDF
\usepackage[warn]{mathtext}
\usepackage[T2A]{fontenc}			% кодировка
\usepackage[utf8]{inputenc}			% кодировка исходного текста
\usepackage[english,russian]{babel}	% локализация и переносы
\usepackage{mathtext} 				% русские буквы в формулах
\usepackage{csvsimple}              % for tabular from csv loading
\usepackage{indentfirst}            % indent after sections
%\usepackage{minipage}

%%% Дополнительная работа с математикой
\usepackage{amsmath,amsfonts,amssymb,amsthm,mathtools} % AMS
\usepackage{icomma} % "Умная" запятая: $0,2$ --- число, $0, 2$ --- перечисление

%%% Номера формул
%\mathtoolsset{showonlyrefs=true} % Показывать номера только у тех формул, на которые есть \eqref{} в тексте.
%\usepackage{leqno} % Немуреация формул слева

%%% Шрифты
\usepackage{euscript}	 % Шрифт Евклид
\usepackage{mathrsfs} % Красивый матшрифт

%%% Свои команды
\DeclareMathOperator{\sgn}{\mathop{sgn}}

%%% Перенос знаков в формулах (по Львовскому)
\newcommand*{\hm}[1]{#1\nobreak\discretionary{}
{\hbox{$\mathsurround=0pt #1$}}{}}

%%% Работа с картинками
\usepackage{graphicx}  % Для вставки рисунков
\graphicspath{{images/}{images2/}}  % папки с картинками
\setlength\fboxsep{3pt} % Отступ рамки \fbox{} от рисунка
\setlength\fboxrule{1pt} % Толщина линий рамки \fbox{}
\usepackage{wrapfig} % Обтекание рисунков и таблиц текстом

%%% Работа с таблицами
\usepackage{array,tabularx,tabulary,booktabs} % Дополнительная работа с таблицами
\usepackage{longtable}  % Длинные таблицы
\usepackage{multirow} % Слияние строк в таблице

%%% Теоремы
\theoremstyle{plain} % Это стиль по умолчанию, его можно не переопределять.
%\newtheorem{theorem}{Теорема}[section]
%\newtheorem{proposition}[theorem]{Утверждение}
 
%\theoremstyle{definition} % "Определение"
%\newtheorem{corollary}{Следствие}[theorem]
%\newtheorem{problem}{Задача}[section]
 
%\theoremstyle{remark} % "Примечание"
%\newtheorem*{nonum}{Решение}

%%% Программирование
\usepackage{etoolbox} % логические операторы

%%% Страница
\usepackage{extsizes} % Возможность сделать 14-й шрифт
\usepackage{geometry} % Простой способ задавать поля
	\geometry{top=25mm}
	\geometry{bottom=35mm}
	\geometry{left=35mm}
	\geometry{right=20mm}
	
%%% Колонтитулы
%\usepackage{fancyhdr}
 	%\pagestyle{fancy}
 	%\renewcommand{\headrulewidth}{0mm}  % Толщина линейки, отчеркивающей верхний колонтитул
 	%\lfoot{Нижний левый}
 	%\rfoot{Нижний правый}
 	%\rhead{Верхний правый}
 	%\chead{Верхний в центре}
 	%\lhead{Верхний левый}
 	% \cfoot{Нижний в центре} % По умолчанию здесь номер страницы
 	
%%% Интерлиньяж
%\usepackage{setspace}
%\onehalfspacing % Интерлиньяж 1.5
%\doublespacing % Интерлиньяж 2
%\singlespacing % Интерлиньяж 1

%%% Гиперссылки
\usepackage{hyperref}
\usepackage[usenames,dvipsnames,svgnames,table,rgb]{xcolor}
\hypersetup{				% Гиперссылки
    unicode=true,           % русские буквы в раздела PDF
    pdftitle={Заголовок},   % Заголовок
    pdfauthor={Автор},      % Автор
    pdfsubject={Тема},      % Тема
    pdfcreator={Создатель}, % Создатель
    pdfproducer={Производитель}, % Производитель
    pdfkeywords={keyword1} {key2} {key3}, % Ключевые слова
    colorlinks=true,       	% false: ссылки в рамках; true: цветные ссылки
    linkcolor=red,          % внутренние ссылки
    citecolor=green,        % на библиографию
    filecolor=magenta,      % на файлы
    urlcolor=cyan           % на URL
}

%%% Другие пакеты
\usepackage{lastpage} % Узнать, сколько всего страниц в документе.
\usepackage{soul} % Модификаторы начертания
\usepackage{csquotes} % Еще инструменты для ссылок
%\usepackage[style=authoryear,maxcitenames=2,backend=biber,sorting=nty]{biblatex}
\usepackage{multicol} % Несколько колонок
\usepackage{multirow} % Несколько строк

%%% Шрифты
%\renewcommand{\familydefault}{\sfdefault} % Начертание шрифта


%%% Работа с библиографией
%\usepackage{cite} % Работа с библиографией
%\usepackage[superscript]{cite} % Ссылки в верхних индексах
%\usepackage[nocompress]{cite} % 
%\usepackage{csquotes} % Еще инструменты для ссылок


%%% Tikz
\usepackage{tikz} % Работа с графикой
\usepackage{pgfplots} % Работа с pgf
\usepackage{pgfplotstable}
\usepackage{upgreek}

%%% Дополнительные пакеты для tikz
%\usepgfplotslibrary{dateplot} % Возможность подписания дат
\pgfplotsset{compat=1.5}

\begin{document}
	\newcommand{\HRule}{\rule{\linewidth}{0.7mm}} % Defines a new command for the horizontal lines, change thickness here
	
	\begin{center}
		\large\textbf{Московский Физико-Технический Институт}\\ % Name of your university/college
		\large\textbf{(государственный университет)}
	
		\vfill
		
		\Large Лабораторная работа по курсу общей физики № *labnum*\\[0.5cm] % Preambule of your document title
		
		
		\HRule
		\\[0.4cm]
		{ \huge \bfseries *name of your labwork*}% Title of your document
		\\[0.4cm] 
		\HRule
		\\[0.5cm]
		
		\ \\
	\textbf{\large Автор:} \\	
	\large *your name* *groupname*\\ % Your name and something more, your group num for example
		\vfill
		\hspace*{-0.8 cm}\includegraphics[width=100 pt]{frkt_logo}\\ % logo of your  company/university/college
		\large Долгопрудный, 2021 % location and year
	\end{center}

\newpage
\setcounter{page}{2}
\fancyfoot[c]{\thepage}
\fancyhead[L] {Работа № *labnum*} % some information in page header
\fancyhead[R]{}
	
	\section*{Задание 1.}
	
	Соберем на макетной плате схема интегрирующей цепи с параметрами $R = 100$ Ом, $C = 1,05$ мкФ, тогда постоянная времени для этой цепи $\tau = R C = 105$ мкс.
	
	Экспериментально определим верхнюю граничную частоту $\nu_0$, подбирая $\nu_0$ таким образом, чтобы амплитуда выходного сигнала составила $70 \%$ от амплитуды входного сигнала. $U_{\text{вых}} = 0,7 U_{\text{вх}}$.
	
	Теоретическое значение $\nu_{\text{0теор}} = \frac{1}{2 \pi C R} = 1,516$ кГц. Экспериментальное значение $\nu_{\text{0экс}} = 1,51$ кГц.
	
	%---------------------------------------------------------------------------------------------------------------------------------------------------------------
	
	Снимем зависимость коэффициента передачи $K(\nu) = \frac{U_{\text{вых}}}{U_{\text{вх}}}$ от частоты $\nu = 2^n \nu_0$ в пределах $n \in [-2, 4]$. Построим график $K(\nu)$, а так же граф Боде -- $20 \lg K$ от$n = \log_2 (\nu / \nu_0)$.

	\begin{table}[h!]
		\begin{center}
			\begin{tabular}{|c|c|c|c|}
				\hline
				$n$& $\nu$, кГц & $K(\nu)$ &  $20 \lg K$      \\ \hline
				-2 & 0,377      & 0,96     & -0,35  \\ \hline
				-1 & 0,755      & 0,9      & -0,92  \\ \hline
				0  & 1,51       & 0,73     & -2,73  \\ \hline
				1  & 3,02       & 0,55     & -5,19  \\ \hline
				2  & 6          & 0,32     & -9,90  \\ \hline
				3  & 12         & 0,16     & -15,92 \\ \hline
				4  & 24,1       & 0,088    & -21,11 \\ \hline
			\end{tabular}
		\end{center}
		\caption{}
	\end{table}

	% график 1
	% график 2
	% картинка колебаний
	
	%---------------------------------------------------------------------------------------------------------------------------------------------------------------
	
	По осцилограмме прямоугольных сигналов оценим постоянную времени $\tau$,\\ измерив время нарастания фронта импульса от нуля до уровня $1 - 1/e \approx 0,63$.
	
	Экспериментально получено $\tau = 111$ мкс. Тогда $\nu_0 = \frac{1}{2 \pi \tau} \approx 1434$ Гц.
	
	% картинка прямоугольных сигналов. 
	
	%---------------------------------------------------------------------------------------------------------------------------------------------------------------
	
	Соберем на макетной плате схема интегрирующей цепи. Экспериментально определим нижнюю граничную частоту $\nu_0$.
	
	Теоретическое значение $\nu_{\text{0теор}} = \frac{1}{2 \pi C R} = 1,516$ кГц. Экспериментальное значение $\nu_{\text{0экс}} = 1,9$ кГц.
	
	%---------------------------------------------------------------------------------------------------------------------------------------------------------------
	
	Снимем зависимость коэффициента передачи $K(\nu) = \frac{U_{\text{вых}}}{U_{\text{вх}}}$ от частоты $\nu = 2^n \nu_0$ в пределах $n \in [-4, 2]$. Построим график $K(\nu)$, а так же граф Боде -- $20 \lg K$ от$n = \log_2 (\nu / \nu_0)$.
	
	
	\begin{table}[h!]
		\begin{center}
			\begin{tabular}{|c|c|c|c|}
				\hline
				n  & $\nu$, Гц & $K(\nu)$ & $20 \lg K$  \\ \hline
				-4 & 118,7     & 0,07     & -23,10 \\ \hline
				-3 & 237,5     & 0,11     & -19,17 \\ \hline
				-2 & 475       & 0,246    & -12,18 \\ \hline
				-1 & 950       & 0,43     & -7,33  \\ \hline
				0  & 1900      & 0,69     & -3,22  \\ \hline
				1  & 3800      & 0,87     & -1,21  \\ \hline
				2  & 7600      & 0,95     & -0,45  \\ \hline
			\end{tabular}
		\end{center}
		\caption{]
	\end{table}
	
	% график 1
	% график 2
	% картинка колебаний
	
	По осцилограмме прямоугольных сигналов оценим постоянную времени $\tau$,\\ измерив время спада вершины импульса от нуля до уровня $1/e \approx 0,37$.
	
	Экспериментально получено $\tau = 628,4$ мкс. Тогда $\nu_0 = \frac{1}{2 \pi \tau} \approx 1591$ Гц.
	
	% картинка прямоугольных сигналов. 

\end{document}